\renewcommand{\Anno}{2012}

\renewcommand{\NumeroProblema}{2}

\renewcommand{\Testo}{
La figura sottostante rappresenta i 7 tetramini possibili (i pezzi del classico gioco del Tetris).
\begin{center}
	\psset{unit=3ex}
	\pspicture(0,0)(4,2)
	\psline[linewidth=1.5\pslinewidth](0,0)(4,0)(4,1)(0,1)(0,0)
	\psline[linestyle=dashed,linewidth=0.7\pslinewidth](1,0)(1,1)
	\psline[linestyle=dashed,linewidth=0.7\pslinewidth](2,0)(2,1)
	\psline[linestyle=dashed,linewidth=0.7\pslinewidth](3,0)(3,1)
	\endpspicture
	\hfill
	\pspicture(0,0)(2,2)
	\psline[linewidth=1.5\pslinewidth](0,0)(2,0)(2,2)(0,2)(0,0)
	\psline[linestyle=dashed,linewidth=0.7\pslinewidth](1,0)(1,2)
	\psline[linestyle=dashed,linewidth=0.7\pslinewidth](0,1)(2,1)
	\endpspicture
	\hfill
	\pspicture(0,0)(3,2)
	\psline[linewidth=1.5\pslinewidth](0,0)(3,0)(3,1)(2,1)(2,2)(1,2)(1,1)(0,1)(0,0)
	\psline[linestyle=dashed,linewidth=0.7\pslinewidth](1,0)(1,1)(2,1)(2,0)
	\endpspicture
	\hfill
	\pspicture(0,0)(3,2)
	\psline[linewidth=1.5\pslinewidth](0,0)(3,0)(3,2)(2,2)(2,1)(0,1)(0,0)
	\psline[linestyle=dashed,linewidth=0.7\pslinewidth](1,0)(1,1)
	\psline[linestyle=dashed,linewidth=0.7\pslinewidth](2,0)(2,1)(3,1)
	\endpspicture
	\hfill
	\pspicture(0,0)(3,2)
	\psline[linewidth=1.5\pslinewidth](0,0)(3,0)(3,1)(1,1)(1,2)(0,2)(0,0)
	\psline[linestyle=dashed,linewidth=0.7\pslinewidth](2,0)(2,1)
	\psline[linestyle=dashed,linewidth=0.7\pslinewidth](1,0)(1,1)(0,1)
	\endpspicture
	\hfill
	\pspicture(0,0)(3,2)
	\psline[linewidth=1.5\pslinewidth](0,0)(2,0)(2,1)(3,1)(3,2)(1,2)(1,1)(0,1)(0,0)
	\psline[linestyle=dashed,linewidth=0.7\pslinewidth](1,0)(1,1)(2,1)(2,2)
	\endpspicture
	\hfill
	\pspicture(0,0)(3,2)
	\psline[linewidth=1.5\pslinewidth](3,0)(3,1)(2,1)(2,2)(0,2)(0,1)(1,1)(1,0)(3,0)
	\psline[linestyle=dashed,linewidth=0.7\pslinewidth](1,2)(1,1)(2,1)(2,0)
	\endpspicture
\end{center}

\begin{enumerate}
	\renewcommand{\labelenumi}{(\alph{enumi})}
	\item Determinare per quale o quali dei pezzi precedenti vale la seguente proprietà: ``scartando quel pezzo, è possibile utilizzare i restanti 6 pezzi una ed una sola volta in modo da comporre un rettangolo $6\times 4$''.

	\item  Stessa domanda, nel caso di rettangoli $8\times 3$.
\end{enumerate}

{\small Nota: si intende che nel comporre il rettangolo è possibile traslare, ruotare ed accostare i pezzi a piacimento, ma non è possibile tagliarli o ruotarli nello spazio uscendo dal piano (così, per esempio, non è possibile trasformare l'uno
nell'altro i 2 pezzi a forma di L).}
}

\renewcommand{\Risposta}{
In entrambi i casi la costruzione è possibile se e solo se si
scarta il pezzo a T.}

\renewcommand{\Aiutino}{
Esibire esplicitamente una costruzione dei 2 rettangoli ottenuta utilizzando tutti i pezzi, tranne quello a T.

Supponiamo poi di colorare i rettangoli a scacchiera. Quando andiamo a disporre i pezzi, quanti quadrettini occupano dell'uno e dell'altro colore? Cosa ha di speciale il pezzo a T?
}

\renewcommand{\SoluzioneUfficiale}{\TitoloSoluzione 
È possibile costruire i 2 rettangoli in questione se e solo se si scarta il pezzo a T.

Per mostrare che scartando il pezzo a T la costruzione è possibile, basta esibire 2 esempi espliciti di costruzione, come nella figura seguente.
\begin{center}
	\psset{unit=3ex}
	\hfill
	\pspicture(0,0)(6,4)
	\multips(1,0)(1,0){5}{
	\psline[linestyle=dashed,linewidth=0.7\pslinewidth](0,0)(0,4)
	}
	\multips(0,1)(0,1){3}{
	\psline[linestyle=dashed,linewidth=0.7\pslinewidth](0,0)(6,0)
	}
	\psline[linewidth=1.5\pslinewidth](0,0)(6,0)(6,4)(0,4)(0,0)
	\psline[linewidth=1.5\pslinewidth](1,0)(1,1)(5,1)(5,0)
	\psline[linewidth=1.5\pslinewidth](0,2)(1,2)(1,3)(5,3)(5,2)(6,2)
	\psline[linewidth=1.5\pslinewidth](2,1)(2,3)
	\psline[linewidth=1.5\pslinewidth](4,1)(4,3)
	\psline[linewidth=1.5\pslinewidth](3,4)(3,3)
	\endpspicture
	\hfill
	\pspicture(0,0)(8,3)
	\psline[linewidth=1.5\pslinewidth](0,0)(8,0)(8,3)(0,3)(0,0)
	\multips(1,0)(1,0){7}{
	\psline[linestyle=dashed,linewidth=0.7\pslinewidth](0,0)(0,3)
	}
	\multips(0,1)(0,1){2}{
	\psline[linestyle=dashed,linewidth=0.7\pslinewidth](0,0)(8,0)
	}
	\psline[linewidth=1.5\pslinewidth](2,0)(2,1)(6,1)(6,0)
	\psline[linewidth=1.5\pslinewidth](3,1)(3,3)
	\psline[linewidth=1.5\pslinewidth](5,1)(5,3)
	\psline[linewidth=1.5\pslinewidth](0,1)(1,1)(1,2)(3,2)
	\psline[linewidth=1.5\pslinewidth](8,1)(7,1)(7,2)(5,2)
	\endpspicture
	\hfill\mbox{}
\end{center}

Per mostrare che la costruzione non è possibile se si scartano altri pezzi, ricorriamo ad una tecnica di colorazione. Supponiamo di colorare a scacchiera, nel modo classico, i rettangoli finali. Ogni rettangolo conterrà un ugual numero di quadratini dei 2 colori. È inoltre abbastanza semplice verificare che i tetramini diversi da quello a T, comunque disposti nel rettangolo finale, occuperanno sempre 2 caselle di un colore e 2 caselle dell'altro, mentre il tetramino a~T occuperà sempre 3 caselle di un colore ed una sola dell'altro. La figura sottostante mostra solo alcune delle possibili disposizioni dei pezzi, ma la sostanza non cambia se anche i pezzi vengono traslati o ruotati.
\begin{center}
	\psset{unit=3ex}
	\pspicture(0,0)(4,2)
	\psframe*[linecolor=green](0,0)(1,1)
	\psframe*[linecolor=green](2,0)(3,1)
	\psline[linewidth=1.5\pslinewidth](0,0)(4,0)(4,1)(0,1)(0,0)
	\psline[linestyle=dashed,linewidth=0.7\pslinewidth](1,0)(1,1)
	\psline[linestyle=dashed,linewidth=0.7\pslinewidth](2,0)(2,1)
	\psline[linestyle=dashed,linewidth=0.7\pslinewidth](3,0)(3,1)
	\endpspicture
	\hfill
	\pspicture(0,0)(2,2)
	\psframe*[linecolor=green](0,0)(1,1)
	\psframe*[linecolor=green](1,1)(2,2)
	\psline[linewidth=1.5\pslinewidth](0,0)(2,0)(2,2)(0,2)(0,0)
	\psline[linestyle=dashed,linewidth=0.7\pslinewidth](1,0)(1,2)
	\psline[linestyle=dashed,linewidth=0.7\pslinewidth](0,1)(2,1)
	\endpspicture
	\hfill
	\pspicture(0,0)(3,2)
	\psframe*[linecolor=green](0,0)(1,1)
	\psframe*[linecolor=green](2,0)(3,1)
	\psframe*[linecolor=green](1,1)(2,2)
	\psline[linewidth=1.5\pslinewidth](0,0)(3,0)(3,1)(2,1)(2,2)(1,2)(1,1)(0,1)(0,0)
	\psline[linestyle=dashed,linewidth=0.7\pslinewidth](1,0)(1,1)(2,1)(2,0)
	\endpspicture
	\hfill
	\pspicture(0,0)(3,2)
	\psframe*[linecolor=green](0,0)(1,1)
	\psframe*[linecolor=green](2,0)(3,1)
	\psline[linewidth=1.5\pslinewidth](0,0)(3,0)(3,2)(2,2)(2,1)(0,1)(0,0)
	\psline[linestyle=dashed,linewidth=0.7\pslinewidth](1,0)(1,1)
	\psline[linestyle=dashed,linewidth=0.7\pslinewidth](2,0)(2,1)(3,1)
	\endpspicture
	\hfill
	\pspicture(0,0)(3,2)
	\psframe*[linecolor=green](0,0)(1,1)
	\psframe*[linecolor=green](2,0)(3,1)
	\psline[linewidth=1.5\pslinewidth](0,0)(3,0)(3,1)(1,1)(1,2)(0,2)(0,0)
	\psline[linestyle=dashed,linewidth=0.7\pslinewidth](2,0)(2,1)
	\psline[linestyle=dashed,linewidth=0.7\pslinewidth](1,0)(1,1)(0,1)
	\endpspicture
	\hfill
	\pspicture(0,0)(3,2)
	\psframe*[linecolor=green](0,0)(1,1)
	\psframe*[linecolor=green](1,1)(2,2)
	\psline[linewidth=1.5\pslinewidth](0,0)(2,0)(2,1)(3,1)(3,2)(1,2)(1,1)(0,1)(0,0)
	\psline[linestyle=dashed,linewidth=0.7\pslinewidth](1,0)(1,1)(2,1)(2,2)
	\endpspicture
	\hfill
	\pspicture(0,0)(3,2)
	\psframe*[linecolor=green](1,1)(2,2)
	\psframe*[linecolor=green](2,0)(3,1)
	\psline[linewidth=1.5\pslinewidth](3,0)(3,1)(2,1)(2,2)(0,2)(0,1)(1,1)(1,0)(3,0)
	\psline[linestyle=dashed,linewidth=0.7\pslinewidth](1,2)(1,1)(2,1)(2,0)
	\endpspicture
\end{center}

Se dunque fosse possibile comporre un rettangolo usando il tetramino a~T più altri 5 tetramini, quel rettangolo finirebbe necessariamente per avere più caselle di un colore che di un altro, il che è assurdo.
}

\renewcommand{\SoluzioniAlternative}{
}

\renewcommand{\Osservazioni}{
La soluzione presentata assume che la costruzione del rettangolo avvenga con i pezzi che ``rispettano la quadrettatura''.  Si può dimostrare che tale ipotesi, oltre ad essere sicuramente intuitiva, è in realtà proprio una condizione necessaria, nel senso che la costruzione non può avvenire altrimenti. Ad esempio si può osservare che i pezzi che occupano i vertici del rettangolo dovranno necessariamente rispettare la quadrettatura, e poi ripetere il ragionamento per i pezzi che occuperanno i vertici della parte rimanente.
}

\renewcommand{\Commenti}{
}